\documentclass[11pt,a4paper]{article}
\usepackage[margin=25mm]{geometry}
\usepackage{amsmath,amssymb,amsthm,hyperref}
\newtheorem{lemma}{Lemma}
\newcommand{\gen}[1]{\langle#1\rangle}
\title{The remaining mixed orbit defect}
\author{Human Compression project: bounded L044 proof component}
\date{}
\begin{document}\maketitle
Retain the accepted unique-crossing setting in a transitive string
C-group $G\leq A_n$, $n\geq12$, with every maximal omission
intransitive. Its two $G_i$-orbits $A,B$ have maximal block systems
$(k_A,m_A)=(k,2)$ and $(k_B,m_B)=(2,m)$, $m\geq3$, and crossing
endpoints $a,b$. Assume $\alpha=0$, $\beta=1$, $\gamma=0$.
Here $M$ is the minimal distinguished subfamily generating the
simultaneous block image, with the accepted actual-orbit deficits.
Then $|M|=m$ and $r=k+m$. Write
$L=\gen{\rho_j:j<i}$ and $R=\gen{\rho_j:j>i}$, so
$G_i=LR$ and $L,R$ commute. We distinguish block images from their
actual permutation lifts throughout.

\begin{lemma}[The block-image marking]
After reversal, $M$ consists of $\tau=\rho_{i-1}$ and the interval
$h_0,\ldots,h_{m-2}=\rho_{i+1},\ldots,\rho_{i+m-1}$.
The image of $\tau$ is pure $C_2$ on the two blocks of $A$;
the other images form a pure maximal anchored $S_m$ marking
on the blocks of $B$. These are image statements; the actual
group $H=\gen{h_0,\ldots,h_{m-2}}$ may have a block kernel.
\end{lemma}
\begin{proof}
The independent block-image family has length $m$, so its group
is the full $C_2\times S_m$ by the paid normal-quotient and
Whiston equality argument. Repeat the block-image part of Lemma~1
of the accepted mixed-block note: a minimal subfamily generating
$S_m$ has $m-1$ members forming a connected extremal string
interval next to $i$. The one remaining generator has trivial
$S_m$ image, by the primitive interval lemma on the opposite
side, or the anchored endpoint-omission centralizer on the same
side. It is the pure $C_2$ element and must have index $i-1$.

This argument does not assume IP in the image. In the next step
of the earlier proof, replace the use of IP by minimality of
$M$ in its simultaneous block image: if the subgroup generated
by the $S_m$ interval contained the pure $C_2$ element, that
remaining generator would be redundant in the block image.
Thus the interval is a character graph over $S_m$. Its tail
fixes the crossing block and generates the natural $S_{m-1}$;
some tail image is odd. Since all tail indices are distant
from $i$, they have trivial binary image. The character is
therefore trivial. This proves the statement entirely in the
block image, without asserting that its lift is faithful.
\end{proof}

The groups $\gen\tau$ and $H$ commute by their actual indices.
Every $M$-orbit in $A$ has even size, and every $M$-orbit in
$B$ meets every one of its $m$ blocks. The single orbit deficit
therefore has precisely two possible locations:
\begin{enumerate}
\item $M$ has one four-point orbit $D$ and $k-2$ paired orbits
in $A$, and two natural $m$-point copies in $B$;
\item $M$ has $k$ paired orbits in $A$ and is transitive on all
$2m$ points of $B$.
\end{enumerate}

\begin{lemma}[The defect on the two-block orbit]
The first alternative is impossible.
\end{lemma}
\begin{proof}
On $B$, the two $M$-orbits meet each block in one point.
The action there is faithfully the diagonal natural $S_m$,
and $\tau$ is trivial there. On every paired $M$-orbit in $A$,
$H$ is trivial because its binary block image is trivial.
On the unique orbit $D$, $\tau$ exchanges the two blocks,
while $H$ preserves them and commutes with $\tau$. Thus the
image of $H$ on $D$ is the diagonal $C_2$ on its two
within-block pairs. Denote its nonidentity element by $z$,
extended by the identity outside $D$.

If $a\notin D$, every right generator fixes $a$, so the
commuting anchor lemma makes the whole right interval trivial
on $A$. This contradicts the nontrivial $H$-action on $D$.
Hence $a\in D$. Only $h_0$ can act nontrivially on $D$:
each later $h_j$ fixes $a$, and a nonidentity element of this
diagonal $C_2$ has no fixed point on $D$.

Consequently the actual $H$ embeds in $\gen z\times S_m$.
Its tail has trivial $z$ coordinate and generates the natural
$S_{m-1}$ on the nonanchor blocks of $B$. The first member
has nontrivial $z$ coordinate. If $z\notin H$, this subgroup
would be a character graph over $S_m$, impossible because the
tail contains an odd image but has trivial $z$ coordinate.
Thus the pure kernel element $z$ really belongs to $H$.
This is a proved kernel occurrence, not a conclusion drawn
from block-image notation.

The entire left interval commutes with $H$ and hence preserves
the support $D$ of $z$. The right generators in $H$ preserve
$D$ by construction. Every right generator beyond $H$ fixes
$a$ and commutes with $h_0$ and $\tau$: the first possible
such index is $i+m$, at distance $m-1\geq2$ from $h_0$.
The group $\gen{\tau,h_0}$ is transitive on $D$, so those
further generators fix $D$ pointwise. Thus $G_i$ preserves
$D$, contradicting transitivity on $A$ unless $k=2$.
When $k=2$, both original systems have two-point blocks;
this is already excluded in the present rank-failure range
by the complete accepted two-point-block lemma. Its range
applies because $n\geq12$ then forces $m\geq4$.
\end{proof}

\begin{lemma}[The defect on the two-point-block orbit]
The second alternative is impossible.
\end{lemma}
\begin{proof}
The group $H$ fixes $A$ pointwise: all paired $M$-orbits
project bijectively to the two blocks, and $H$ has trivial
binary image. Every further right generator fixes $a$ too.
The commuting anchor lemma therefore makes the entire right
interval trivial on $A$. The left interval is transitive on
$A$, though its action on $B$ need not be trivial.

The block image of $H$ on $B$ is transitive $S_m$, so its
point action is either transitive on all $2m$ points, or
has two $m$-point orbits, each a faithful natural $S_m$
copy. In either case an element centralizing $H$ and fixing
$b$ is the identity on $B$: use transitivity in the first
case; in the second, fixing one point fixes that whole copy,
and the other natural copy has trivial internal centralizer.
Thus every left distinguished generator except $\tau$ is
trivial on $B$, because it centralizes $H$ and fixes $b$.
We do not impose this on $\tau$, which may perform the
uniform within-block flip on $B$.

There is no right generator beyond the $H$ interval.
Any generator more distant than its one outer neighbor
centralizes $H$ and fixes $b$, hence is trivial on $B$
and already trivial on $A$. The one possible outer neighbor
$v=\rho_{i+m}$ centralizes the far-end omission of $H$.
Its block image fixes the crossing block and is trivial by
the anchored symmetric centralizer lemma. It fixes the
crossing two-point block pointwise, and commutation propagates
this fixation to at least $m-1$ blocks. It can therefore
move at most the remaining two-point block, by a single
transposition. Since it is trivial on $A$, evenness excludes
that as well. The entire right interval is exactly $H$,
of rank $m-1$.

The left interval consequently has rank $k$, since
$r=k+m$. It ends in $x=\rho_{i-2}$ and $\tau$, with $x$
present to connect the paired anchor to another row.
Its preceding $k-2$ generators are trivial on $B$, commute
with $\tau$, and fix $a$. They act faithfully and diagonally
on the $k-1$ nonanchor rows of $A$, so their independent
marking generates the full $S_{k-1}$ by Whiston.
The generator $x$ is also trivial on $B$.

For $k\geq5$, the accepted maximal symmetric-string endpoint
classification excludes this equality. For completeness, in
an ordinary marking the near-end omission is $S_{k-2}$,
whose centralizer on a block is the transposition of the
anchor and one unused row. Fixation of $a$ makes $x$
trivial on its block and a single odd transposition on the
other, impossible. The degree-four twisted marking
($k=5$) and outer degree-six marking ($k=7$) can instead
have a transitive endpoint omission on all nonanchor rows.
Any permutation centralizing it preserves its unique fixed
point, so $x$ fixes the opposite anchor, again impossible.
These are exactly the previously accepted small/outer
endpoint cases; no new small-group classification is used.

For $k=2$, the same argument is immediate on two-point
blocks: $x$ is one odd transposition. For $k=3$, label
the rows $0,1,2$ with $a=0$ and their mates $0',1',2'$.
The sole preceding generator is
$p=(1\ 2)(1'\ 2')$. Since $x$ fixes $0$, moves $0'$
to another row and is even, it is
\[
        x=(1\ 2)(0'\ j'),\qquad j\in\{1,2\}.
\]
The group $\gen{x,\tau}$ is transitive on the six points:
contract the three $\tau$ pairs, and its row edges are
$1$--$2$ and $0$--$j$, a spanning tree. At $k=4$ the
complete accepted endpoint calculation gives
\[
 p=(1\ 2)(1'\ 2'),\quad q=(2\ 3)(2'\ 3'),\quad
 x=(1\ 2)(0'\ 3')\ \hbox{or}\ (1'\ 2')(0'\ 3').
\]
Then $\gen{q,x,\tau}$ is transitive on the eight points,
by the accepted row path $0$--$3$--$2$--$1$.

In either small case, omit the far generator $p$ from
$G$. Its remaining generators are transitive on $A$.
They still contain $\gen{\tau,H}=\gen M$, which by this
case's hypothesis is transitive on $B$, whether or not
$H$ alone is transitive there. The generator $\rho_i$
joins the two orbits. Thus this global maximal omission
is transitive, the final contradiction.
\end{proof}

\paragraph{Boundary and attribution.}
This completes the last numerical mixed case in the accepted
exact-deficit reduction. Together with the complete two-point-block,
two-block and mixed beta-zero components, it proves the rank
bound $r\leq\lfloor(n-1)/2\rfloor$ of L044 Proposition~5.7
in the standing setting. The
source framework is Cameron--Fernandes--Leemans--Mixer,
arXiv:1605.09173v1, Section~5.1; the complete paid generation,
anchor and endpoint proofs remain explicit dependencies.
The finer source block diagrams, other main source branches
and exact tables still require their own dispositions.
No full-original recovery, new census, Lean theorem or
question credit is claimed.
\end{document}
